%! Principles of Experimental Design — Unit 1, Lesson 1.6 — Lesson Plan
% PILOT SHAPE (2026-09-06): modeled on AP Statistics lesson 1.4 minus its AP Practice page.
% GRADUAL-RELEASE plan (warm-up / guided notes and practice / spoken debrief / close and START
% THE HOMEWORK). The guided notes carry the release ROW BY ROW in one Main Ideas / Notes table
% (2026-09-12 shape): I do / we do in rows 1-4 (problems 1-13) -> we do the Guided Practice row
% (problems 14-17). THERE IS NO INDEPENDENT PRACTICE SET IN THE NOTES — the homework is the
% individual practice, and the period ends with students starting it. Teacher-facing. Breakable
% boxes are `skillbox` with a \boxguard ahead; the two tabularx boxes are `fixedskillbox`
% (never splits). No group activity, no tier boxes, no exit ticket — the legacy activity/ and
% exit_ticket/ directories were folded in and removed on the 2026-09-14 regeneration.
\documentclass[10pt]{article}
\usepackage{saar-article}
\usepackage{saar-boxes}
\usepackage{graphicx}   % -article does NOT load graphicx
\graphicspath{{images/}}
\newcommand{\TallMath}[1]{$\displaystyle #1\rule[-1.4em]{0pt}{3.2em}$}
\newcommand{\UnitNumberName}{Unit 1: Asking Questions with Data: Study Design \quad}
\newcommand{\LessonNumberName}{Lesson 1.6: Principles of Experimental Design}

\begin{document}
\begin{center}
    {\Large\bfseries \CourseName \\
    {\small\bfseries \UnitNumberName \LessonNumberName}}
\end{center}

\begin{tcolorbox}[colback=frost, colframe=cerulean, boxrule=0.9pt, arc=2mm,
    left=2mm, right=2mm, top=1.5mm, bottom=1.5mm]
    \textbf{Primary Objective:} Students will tell an \emph{experiment} from an observational
    study by asking who built the groups, name the experimental units, the treatment, and the
    control group in a described design, check that design against \emph{comparison,
    randomization, replication,} and \emph{control} and name the principle a changed plan
    breaks, split units into blocks and say who is blinded --- and justify why random
    assignment, not the size of the gap, is what lets a study say \emph{caused}.
    \\[2pt]
    \textbf{Standards:} PS.DC.3a (i--iv), PS.DC.3b \quad$\cdot$\quad AFDA.DA.2c
    \quad$\cdot$\quad PS.DC.1b carried forward in the conclusion items
    \quad$\cdot$\quad answers Lesson 1.5 (PS.DC.2d)
    \\[2pt]
    \textbf{Lesson model:} \emph{Gradual release --- I do, we do, you do --- all of it inside
    the guided notes.} There is no group activity and no exit ticket. The notes name the
    vocabulary as they teach it and deliver the instruction \textbf{row by row in one Main
    Ideas / Notes table} (four rows, problems 1--13: you read the definition, mark up the
    display, and work the first problem of each grid; the class works the rest); the Guided
    Practice row (problems 14--17) is worked together with students holding the pen;
    \textbf{the homework is the individual practice}, and students start it in class while you
    circulate. The whole-class debrief is \emph{spoken}.
    \\[2pt]
    \textbf{Note on scope:} no inference, no significance, no $P$-values. ``Is $20$ points big
    enough to be real?'' is Unit~7. Today the question is only whether the \emph{design} lets
    anyone say \emph{caused} --- say so and move on.
\end{tcolorbox}

\renewcommand{\arraystretch}{1.4}
\boxguard
\begin{skillbox}[Learning Targets \& Key Understandings]{goldbox}
\begin{tabularx}{\linewidth}{@{}X|X@{}}
\textbf{Learning targets (student language)} & \textbf{Key understandings (the ``why'')} \\[2pt]
\hline
\vspace{-6pt}
\begin{itemize}[leftmargin=1.1em, nosep, after=\vspace{2pt}]
  \item I can tell an \textbf{experiment} from an observational study by asking who built the
        groups, and name the \textbf{experimental units} (\textbf{subjects}), the
        \textbf{treatment}, and the \textbf{control group}.
        \hfill \textit{PS.DC.3a.i, a.iii}
  \item I can check a design against \textbf{comparison, randomization, replication,} and
        \textbf{control}, and name the principle a changed plan breaks --- including the
        \textbf{confounding variable} that control keeps out.
        \hfill \textit{PS.DC.3b; AFDA.DA.2c}
  \item I can explain the \textbf{placebo effect}, say who is \textbf{blinded}
        (\textbf{single-} vs.\ \textbf{double-blind}), and split units into the blocks of a
        \textbf{randomized block design}. \hfill \textit{PS.DC.3a.ii, a.iv}
  \item I can justify why \textbf{random assignment} lets an experiment say \emph{caused}
        when watching could only say \emph{associated}. \hfill \textit{PS.DC.3b; 1b}
\end{itemize}
&
\vspace{-6pt}
\begin{itemize}[leftmargin=1.1em, nosep, after=\vspace{2pt}]
  \item \emph{An experiment assigns; an observational study watches.} That one word is the
        whole difference, and it is exactly the thing Lesson~1.5 said was missing.
  \item \textbf{Target misconception:} that a bigger gap is stronger evidence. Riverbend's
        \emph{watched} study gave a $15$-point gap and its \emph{assigned} experiment gives
        $20$; Harbor Point's watched gap was $25$ and its assigned gap is $20$; Millbrook's
        was $48$ and is now $30$. \emph{The design buys the verb --- not the gap.} A
        self-selected gap is inflated by whoever did the selecting.
  \item \textbf{Second misconception:} that a control group is a wasted group, and that a
        group which gets \emph{something else} is not a control. Without it there is nothing
        to compare against, and comparison is principle one (notes problem~4).
  \item \textbf{Randomization} makes the groups alike \emph{before} the treatment;
        \textbf{control} keeps them alike \emph{during} it. A \textbf{confounding variable} is
        1.5's lurking variable that got \emph{inside} the experiment.
  \item Blinding does not change who was treated --- it removes the part of the response that
        comes from \emph{knowing} you were picked.
  \item \textbf{Blocking} does not replace chance; it applies chance inside groups of similar
        units so a variable you already suspect cannot land unevenly.
\end{itemize}
\end{tabularx}
\end{skillbox}

\boxguard[24]
\begin{skillbox}[{Vocabulary, Concepts, \& Theorems --- taught directly in the guided notes}]{greenbox}
\small
\renewcommand{\arraystretch}{1.35}
\begin{tabularx}{\linewidth}{@{} >{\bfseries}p{3.1cm} X @{}}
Experiment & A study in which the researcher assigns the units to groups and imposes a
       treatment on at least one group. \\
Experimental unit & The individual or object a treatment is applied to. When the units are
       people they are \textbf{subjects}. \\
Treatment & The condition the researcher applies to a group. \\
Control group & The group that does not get the treatment --- the baseline the treatment
       group is compared against. It is not a wasted group. \\
Random assignment & Using chance --- a generator, a coin, drawn slips --- to decide which
       group each unit goes in, so the groups start out alike. \\
Four principles & \textbf{Comparison} (at least two groups), \textbf{randomization} (chance
       decides), \textbf{replication} (many units per group), \textbf{control} (everything
       else the same). \\
Confounding variable & A variable that differs between the groups \emph{and} affects the
       response, so the two explanations cannot be told apart. \\
Placebo effect & A subject responds to the \emph{idea} of being treated. A \textbf{placebo}
       is a fake treatment built to look real. \\
Blinding & Keeping someone from knowing which group a subject is in.
       \textbf{Single-blind:} the subjects, or the measurers. \textbf{Double-blind:} both. \\
Completely randomized & Every unit is assigned purely by chance (Riverbend's $60/60$). \\
Randomized block & Units are sorted into \textbf{blocks} of similar units first, then
       randomized \emph{within} each block. \textbf{Matched pairs} is the block of two. \\
\end{tabularx}
\end{skillbox}

% fixedskillbox never splits, so the phase table stays intact on one page.
% THE PHASE TABLE IS A CONTRACT: 5 / 35 / 10 / 10 — the I do is ~20 of the 35 and the Guided
% Practice ~15; the debrief is spoken; the homework start is the second formative read.
\begin{fixedskillbox}[Lesson at a Glance --- \MeetingLength]{frost}
\renewcommand{\arraystretch}{1.25}
\begin{tabularx}{\linewidth}{@{}l c X X@{}}
\textbf{Phase} & \textbf{Min} & \textbf{Students} & \textbf{Teacher} \\
\hline
Warm-Up & 5 & The two Riverbend group rates, the $20$-point gap, and last quarter's $15$-point
  gap the other way --- silent & Debrief item~3 aloud; write \emph{assigned: $75$ vs.\ $55$
  \quad watched: $65$ vs.\ $80$} on the board and leave it up all period \\
Guided Notes \& Practice & 35 & Fill the vocabulary box as each term is named; work rows 1--4
  of the table (problems 1--13) with the pen in their hand; then the Guided Practice row,
  \emph{Harbor Point Swim Times} (14--17), together & \textbf{I do / we do, row by row}: read
  the definition, mark up the display, work problems 1, 6, 8, 10; the class works the rest
  (20) $\to$ \textbf{we do} Guided Practice 14--17 (15) \\
Debrief & 10 & Pens down, papers up --- answer aloud, nothing to fill in & Put the two
  Riverbend studies back up and take the ``the bigger gap wins'' reflex, then the
  ``a control group is wasted'' reflex; cold check on Bayside \\
Close \& Assign & 10 & \textbf{Start the homework in class}, alone and silent & Announce the
  homework and its form (packet or Desmos), circulate for the second formative read, preview
  Lesson~1.7 \\
\end{tabularx}
\end{fixedskillbox}

\begin{fixedskillbox}[Warm-Up --- Activate Prior Knowledge (5 min)]{frost}
\begin{tabularx}{\linewidth}{@{}X X@{}}
\begin{minipage}[t]{\linewidth}\vspace{0pt}
    \textbf{The three items and what each seeds}
    \begin{itemize}[leftmargin=1.1em, itemsep=1pt, topsep=2pt]
      \item \textbf{1.} $45$ of $60$ as a percent ($75\%$). \textbf{Seeds:} the treatment
            group's rate --- it is the number in row~1's display and in problem~6's
            \emph{replication} check.
      \item \textbf{2.} $33$ of $60$ ($55\%$). \textbf{Seeds:} the \emph{control} group's
            rate. Note aloud that both denominators are $60$ --- \emph{because a generator
            built the groups}. That is row~1.
      \item \textbf{3(a).} $75 - 55 = 20$ percentage points. \textbf{Seeds:} the gap row~4
            argues about.
      \item \textbf{3(b).} Last quarter's \emph{watched} result: non-users higher, by $15$
            points. \textbf{Seeds:} the hook --- both gaps are on the board before anyone
            votes, and problem~12 asks whether gap size was ever the reason.
    \end{itemize}
\end{minipage}
&
\begin{minipage}[t]{\linewidth}\vspace{0pt}
    \textbf{Running it}
    \begin{itemize}[leftmargin=1.1em, itemsep=1pt, topsep=2pt]
      \item \textbf{Debrief 3(b) aloud} and write both studies on the board ---
            \emph{assigned: $75$ vs.\ $55$ \quad watched: $65$ vs.\ $80$}. Leave it up all
            period; the hook, row~4, and the debrief all point back at it.
      \item 3(b) is the tell. A student who says the Study Center is harmful has read the
            watched study as proof and needs row~4 slowly. A student who asks \emph{which one
            is right} has already found today's question.
      \item Items 1--2 are one division each --- say so, so nobody thinks they are being
            retested. Do not let them expand.
      \item No vocabulary here on purpose. The words \emph{treatment} and \emph{control
            group} arrive in row~1, hung on these two numbers.
    \end{itemize}
\end{minipage}
\end{tabularx}
\end{fixedskillbox}

\boxguard
\begin{skillbox}[Hook --- before anyone sits down]{frost}
The hook is on \textbf{page 1 of the notes}, under the vocabulary box, and on the board:
\textit{``Last quarter the counselor watched: $65\%$ of Study Center users passed every class
against $80\%$ of the students who never went --- the Study Center looked harmful. This
quarter she assigned: $120$ volunteers split by a generator, $75\%$ against $55\%$ --- now it
looks helpful.''} \textbf{Ask: same school, same Study Center, opposite conclusions --- which
study should she believe?} Students circle \emph{last quarter's}, \emph{this quarter's}, or
\emph{neither} and write one line of reason \emph{before} anyone speaks; then take the hand
vote and write the split on the board. Most rooms pick this quarter's --- and most of them
pick it for the wrong reason (``it is newer,'' ``the gap is bigger'').
\textbf{Do not settle it.} Say only that the answer is problem~12, in the \emph{Caused} row,
and that they will vote again there. The reason matters more than the verdict, and you cannot
show that until they have committed to one.
\end{skillbox}

\boxguard
\begin{skillbox}[Guided Notes \& Practice --- I do, then we do (35 min)]{frost}
\textbf{This is the whole lesson.} There is no group activity, and there is no independent
practice set on this page: \textbf{the homework is the individual practice}, started in class.
The notes are \textbf{one Main Ideas / Notes table}: four instruction rows (problems 1--13),
then the Guided Practice row (14--17). \textbf{The rhythm in every row is the same}: read the
definition aloud, mark up the display, work the first problem of the grid yourself, then hand
the rest to the room with the pen in their hand. The vocabulary box on page~1 is filled
\emph{as each term is named}, never in advance. The only blanks are the four cells of the
Cedar Ridge table (problem~5); every other answer is written in a problem's own space.
\begin{multicols}{2}
\textbf{I do / we do, row by row (about 20 min).}

\emph{Row 1, Experiment (problems 1--4).} Open on the two lines already on the board. Read the
definition, then name the parts of the flow diagram --- the $120$ volunteers are the
\emph{subjects}, twice-a-week attendance is the \emph{treatment}, the other $60$ are the
\emph{control group}. Work problem~1 yourself; the class does 2--4. \textbf{Problem~4 is the
trap}: a new tutoring app against the old worksheet packet. Take a hand count before anyone
writes --- half the room will say there is no control group because \emph{everybody gets
something}. The packet group \emph{is} the control: it is the baseline. Five minutes.

\emph{Row 2, Principles (problems 5--7).} Give the four cards from the display, then take the
four Cedar Ridge plans of problem~5 from the class --- one altered plan per principle, in
order. Work problem~6 yourself on Riverbend's own design so the four words attach to numbers
they already have. Problem~7 is the confounding variable (the new bins) and it is the row's
real content --- do not let it get cut. Six minutes.

\emph{Row 3, Blinding \& Blocking (problems 8--9).} Read the two refinements and mark the
two-panel display. Work problem~8 yourself (the placebo effect, and how Room~B removes it);
the class does problem~9, which is the block arithmetic --- $80 \to 40/40$, $40 \to 20/20$,
$60$ per group --- plus the reason to block at all. Four minutes; this row moves.

\emph{Row 4, Caused (problems 10--13) --- the crux.} Put the two studies side by side on the
display and work problem~10 live (where the struggling students went this quarter). The class
does 11, the conclusion sentence with the verb the design earns. Then \textbf{re-take the hook
vote at problem~12 before you say anything} --- and take the classmate's reason in the problem
(``$20$ points beats $15$'') head-on. \textbf{Problem~13 is the crux}: a \emph{watched} study
with a $35$-point gap, bigger than the experiment's $20$. It still earns only
\emph{association}. Land the printed sentence word by word. Five minutes.

\columnbreak
\textbf{We do (about 15 min).} The Guided Practice row, \emph{Harbor Point Swim Times}
(problems 14--17) --- the same four moves the homework will ask alone, on a pool instead of a
school. \textbf{Students hold the pen; you hold the questions.} Problem~14 goes on them cold
(units, treatment, control, who decided). Problem~15 is warm-up arithmetic on new numbers
($27/45 = 60\%$, $18/45 = 40\%$, a $20$-point gap). Problem~16 is the four principles plus the
honest limit of blinding --- the swimmers cannot be blinded, the sleep-log scorers can.
\textbf{Problem~17 is the crux transferred}: watching gave $25$ points, assigning gives $20$.
Take a hand count on \emph{which is stronger} before anyone writes.

Circulate during Guided Practice and demand the reason, not the word:
\begin{itemize}[leftmargin=1.1em, nosep]
  \item \textbf{Problem 14:} ``Who decided which slot you swam in? Then what is the $7$~p.m.\
        group called?''
  \item \textbf{Problem 15:} ``Both groups have $45$. Why is that not an accident?''
  \item \textbf{Problem 16:} ``You wrote \emph{double-blind}. Can a swimmer not know what time
        she swam?''
  \item \textbf{Problem 17:} ``Last lesson $75\%$ of the morning swimmers were retired. Where
        are the retired members this time?'' (Split across both groups by chance.)
\end{itemize}
While circulating, pick the papers to put up in the debrief --- one problem~17 that said
\emph{the 20, because chance built the groups}, and one that picked the $25$. \textbf{Your
formative read comes twice}: here, and again in the last ten minutes as students start the
homework.
\end{multicols}
\end{skillbox}

\boxguard[24]
\begin{skillbox}[Debrief --- whole class, spoken (10 min)]{frost}
Pens down, papers up. \textbf{This debrief is spoken} --- there is no box at the end of the
notes to fill in, so it lives entirely on the board and in the room.
\begin{multicols}{2}
\begin{enumerate}[leftmargin=1.3em, nosep, itemsep=3pt]
  \item \textbf{Put the two Riverbend studies back up} and make a student say the sentence:
        the experiment may say \emph{caused} because chance built the groups --- not because
        $20$ beats $15$. Then state it as PS.DC.3b: a well-designed experiment is what detects
        a cause-and-effect relationship.
  \item Put up the two problem~17s you picked. The difference between them is not arithmetic
        --- it is whether the student still believes the bigger gap wins. Read both aloud.
  \item Put problem~4 back up and take the \emph{a-control-group-is-wasted} reflex:
        \emph{``The packet group got something. Is it still the control group?''}
  \item Take problem~7 last: the reminder group also got new bins. \emph{Name the confounding
        variable} --- and say why the manager cannot tell which one worked.
\end{enumerate}
\columnbreak
\textbf{Check for understanding --- ask it cold, whole class.} \emph{``Bayside Middle School:
$80$ volunteers, a generator sends $40$ to a daily reading block and $40$ to regular homeroom.
Afterward $28$ of the reading-block students improved and $16$ of the homeroom students did. A
teacher says: `Next year, just give the reading block to the students who want it.' Name the
principle that change breaks, and say what could go wrong.''} Hands down, thirty seconds of
think time, then cold-call three students. Correct answer: it breaks \textbf{randomization}
--- the students would choose their own groups, the groups would no longer start alike, and
whatever made them volunteer (a parent, a study habit, a reading level) would be a
\textbf{confounding variable}; the school could no longer say the block \emph{caused}
anything. \emph{(Rates, if asked: $70\%$ against $40\%$, a $30$-point gap.)} \emph{This
replaces the exit ticket.}

\textbf{Formative read.} Sort the three answers as they come: \textbf{(a)} named
randomization and said the groups would stop being alike; \textbf{(b)} said ``that's not fair''
or ``it's biased'' with no principle; \textbf{(c)} said it would be fine, or thought a bigger
gap next year would prove more. Pile (b) is the teachable one and will be large --- give it
the word. \textbf{If pile (c) runs past a third of the room, open Lesson~1.7 by re-running the
\emph{Caused} row (problems 10--13)}, and say so plainly.

\textbf{If the debrief runs short}, cut items 3 and~4 --- item~1 is the one that cannot be
cut. \textbf{Do not let the debrief eat the last ten minutes} --- the homework start is not
optional padding, it is your second formative read.
\end{multicols}
\end{skillbox}

\boxguard
\begin{skillbox}[Homework --- scored, started in class, due after two study halls]{goldbox}
Six exercises on the \textbf{Millbrook Public Library} ($1{,}200$ card holders; $200$
volunteers split $100/100$ by a generator into the Summer Reading Program and a fall waiting
list): \textbf{1} name the units, treatment, control group, and who decided; \textbf{2} the
two rates and the gap ($65\%$, $35\%$, $30$ points); \textbf{3} pool the $200$ and scale
$50\%$ to $1{,}200$ as an \emph{estimate} (\emph{the spiral to Lesson 1.2}); \textbf{4} the
\textbf{contrast pair} --- the same library watched last summer ($80\%$ vs.\ $32\%$) and
assigned this summer, one may claim an \emph{association} and the other \emph{caused it};
\textbf{5} block on age ($120$ adults $\to 60/60$, $80$ children $\to 40/40$) and say why she
blocked; \textbf{6} \emph{the crux head-on} --- last summer's $48$-point gap against this
summer's $30$, and why the \emph{smaller} gap is the stronger evidence. The homework is capped
at \textbf{two pages} (user decision 2026-09-06).

\textbf{Start it in the last ten minutes of the period, in class and alone.} \textbf{Scoring:}
12 points --- 2 per item. \textbf{Item~6 is where the misconception shows}: a student who
answers ``the $48$ is bigger so it is better'' has learned the words and not the rule.
\textbf{Item~4's two rows carry the second check} --- if both rows say \emph{caused it}, the
student never heard row~4.

\textbf{Packet or Desmos --- decided here.} The packet pages are authored and are the default.
Desmos Classroom's study-design coverage is thin: there is nothing that asks a student to name
the principle a changed plan breaks or to explain why a smaller gap is stronger --- the items
this lesson exists for --- so \textbf{1.6 is a packet night}. Say so aloud at Close \& Assign.

\textbf{Preview of Lesson~1.7.} Millbrook could run an experiment. A study of smoking and lung
disease cannot assign anyone to smoke. Do not resolve that here; it is 1.7's opening --- when
an experiment is worth the cost, and when it is impossible or wrong to run one.

\textbf{Connections \& Big Ideas.} \textit{Unit 1 core idea:} a conclusion is only as
trustworthy as the study that produced it. \textit{Standards covered:} PS.DC.3a (i--iv),
PS.DC.3b, AFDA.DA.2c (with PS.DC.1b carried forward in the conclusion items).
\textit{Spirals forward to:} experiments vs.\ observational studies (1.7), the full-cycle
project (1.8), and --- directly --- inference and significance (7.1--7.3, PS.IS.1).
\textit{Reaches back to:} Lesson 1.5's lurking variable and association-not-causation, and
Lesson 1.2's statistic-as-estimate.
\end{skillbox}

\boxguard
\begin{skillbox}[Watch For (while circulating)]{redbox}
\begin{itemize}[leftmargin=1.2em, nosep]
  \item \textbf{``The bigger gap is the better evidence''} (problems 12, 13, 17; HW~6 --- the
        target misconception). Probe: \emph{``Who built the two groups in the study with the
        bigger gap? Then what else could explain it?''}
  \item \textbf{``There is no control group --- everybody got something''} (problem~4, the
        trap; HW~1). Probe: \emph{``Which group is the baseline you compare against?''}
  \item \textbf{Calling an observational study an experiment} (problems 1 and~3; HW~4 row~i):
        random \emph{selection} mistaken for random \emph{assignment}. Probe: \emph{``Did
        anyone decide which group you were in, or only who got asked?''}
  \item \textbf{Confounding named as ``bias''} (problem~7; HW~4). Probe: \emph{``Name the
        variable. Which group had more of it, and does it affect the response?''}
  \item \textbf{Blinding claimed where it is impossible} (problem~16): swimmers ``can be
        double-blinded.'' Probe: \emph{``Can you not know what time you swam?''}
  \item \textbf{Blocking read as skipping chance} (problem~9; HW~5): ``she put the adults in
        the program.'' Probe: \emph{``What happens \emph{inside} each block?''}
  \item \textbf{The verb} (problems 11 and~17; HW~4, HW~6): \emph{caused} written under a
        watched study, or \emph{associated} written under an experiment. Probe: \emph{``Who
        built the groups?''}
\end{itemize}
Cold-call prompts: \emph{``Give me a treatment that is not a pill.''} \quad \emph{``Name
something about this class you would block on before an experiment.''}
\end{skillbox}

\boxguard
\begin{skillbox}[Close \& Assign (10 min)]{goldbox}
\textbf{Students start the homework here, in class, alone and silent --- this is the individual
practice, and the first minutes of it are your best formative read.} Hand the launch first ---
six exercises on paper, scored out of 12, due at the start of the first class after two study
halls --- and point out that the \emph{Keep in Mind} box on the cover is the page to flip back
to. Circulate: \textbf{item~6 is the column that tells you whether the \emph{Caused} row
landed}, and it is on page~2, so put students on 1, 4, and~6 first and let 2, 3, and~5 go home.
Sort what you see into three piles as you walk --- chose the $30$ and said chance built the
groups (ready); chose the $30$ with no reason (ready, needs the language); chose the $48$, or
explained the difference by the librarians trying harder (the misconception survived) --- and
open Lesson~1.7 by re-running problems 10--13 if that third pile ran past a third of the room.
Then close on what changed: \emph{``You walked in able to say two things move together. You
walk out able to say one \textbf{caused} the other --- but only when chance, not a person and
not the volunteers themselves, built the groups. And you know that the study with the biggest
gap is usually the one that earned it least.''} \textbf{Preview:} Lesson~1.7, \emph{Experiments
and Observational Studies Side by Side} --- Millbrook could assign; nobody can assign a person
to smoke for thirty years. Next time: when an experiment is worth the cost, and when it is
impossible or wrong to run one.
\end{skillbox}

% ── Teacher notes — one per component, in packet order (three) ──────────────
% Teacher-only prose lives HERE, never in a _key: a note is the one block in a key with no
% counterpart in the blank. No note for the debrief or the homework start — both are fully
% specified in their own boxes above.
\begin{teachernote}[Warm-Up]
Five minutes, silent, timed. Items 1 and~2 are one division each --- say so, so nobody thinks
they are being retested. The point worth making aloud is that \textbf{both denominators are
$60$}: nobody chose a group, so the groups came out the same size. That is the first crack of
today's door.

\textbf{Item 3(b) is the one to read while collecting.} It reaches back to Lesson~1.5's
watched study of the same Study Center, where the non-users were $15$ points \emph{higher}.
Students who write that down and look uneasy have found the hook themselves. Students who
write ``the Study Center is bad'' have taken the watched study as proof and will need row~4
slowly. Write both studies on the board in a student's own words ---
\emph{assigned: $75$ vs.\ $55$ \quad watched: $65$ vs.\ $80$} --- and \textbf{leave them up all
period}; you will point at that line in the hook, in row~4, and twice in the debrief. Do
\emph{not} supply the words \emph{treatment} and \emph{control group} here --- the warm-up is
deliberately wordless about them.
\end{teachernote}

\begin{teachernote}[Guided Notes \& Practice]
\textbf{This one document is the lesson --- pace it 20 / 15, then 10 for the debrief, then 10
for the homework start.} It is one table, and every row runs the same way: you read the
definition, mark up the display, and work the first problem of the grid (1, 6, 8, 10); the
class works the rest with the pen in their hand. Rows 1 and~2 must move: five minutes on
assigning-versus-watching, six on the four principles.

\textbf{Take the hand count on problem~4 before anyone writes.} A new tutoring app against the
old worksheet packet: half the room will say there is no control group because both sides got
something. If they write first the trap is gone; if they commit to \emph{no control group} and
\emph{then} hear that the control is whatever you compare against, the definition arrives as
the answer to a question they already have.

\textbf{Spend the time in row~4, \emph{Caused}}, which carries the misconception. Row~3
(blinding and blocking) is four minutes and no more --- give the two refinements, work
problem~8, hand problem~9's arithmetic to the room, move. Then slow down: work problem~10 live
(the struggling students are spread across both groups this quarter), let the class do~11, and
\textbf{re-take the hook vote at problem~12 before you say anything}. Take the classmate's
reason in the problem seriously --- ``$20$ points beats $15$'' is what most of the room voted
on. Then \textbf{problem~13} does the work: a \emph{watched} study with a $35$-point gap, and
it still earns only \emph{association}. Do not resolve it by repeating the definition; ask who
built the groups, twice.

\textbf{Guided Practice (14--17) is where the pen changes hands.} \emph{Harbor Point Swim
Times} runs the four moves the homework will ask alone --- name the parts, compute the rates,
check the principles, judge the two gaps --- on the pool the class has seen since 1.2.
Problem~16 carries the honest limit: you \emph{cannot} blind a swimmer to the time of day, and
saying so is the right answer, not a dodge. Problem~17 is the crux transferred. If you only
get through 14, 15, and~17, that is the right trade; problem~16's principle list can be said
aloud.

\textbf{There is no independent practice set on this page, and that is deliberate --- the
homework is the individual practice.} Guided Practice is the last row; the period ends with
students starting the homework in front of you, which is where your second formative read
comes from. Item~6 is the one to watch. If the ``bigger gap wins'' pile is more than a third of
the room, \textbf{open Lesson~1.7 by re-running problems 10--13}, and say so plainly.

\textbf{Debrief (10 min), spoken.} The two Riverbend studies back on the board with a student
saying why the experiment earns the verb; it is the one move that cannot be cut. \textbf{Do not
let the debrief eat the last ten minutes.}
\end{teachernote}

\begin{teachernote}[Homework]
\textbf{Scored, 12 points (2 per item), due at the start of the first class after two study
halls.} The ten minutes at the end of the period are a \emph{start}, not a completion: put
students on 1, 4, and~6 while you can still catch a wrong start, and let 2, 3, and~5 go home.
Item~6 is the one to read first when scoring --- any student who says the $48$-point gap is
better evidence has learned the words and not the rule. Item~4 is the fullest diagnostic in
the set: the same library, two summers, and only one of the two rows may say \emph{caused it};
a student who writes it twice is still sorting by how big the numbers are. Item~5 catches the
other half of blocking --- a student who assigns all $120$ adults to the program has read
blocking as sorting rather than as randomizing \emph{inside} a block. Score items 1--5 for
correctness and item~6 for the \emph{reason}, not the verdict. \textbf{When you score the page,
sort item~6 into three piles} --- chose the $30$ and named random assignment (ready); chose the
$30$ vaguely (ready, needs the language); chose the $48$ or blamed the librarians (the
misconception survived). If more than a third land in the last pile, reopen the two Riverbend
studies in the Lesson~1.7 warm-up debrief rather than letting it ride.

\textbf{Packet is the call for 1.6.} Desmos has nothing that asks a student to name a broken
principle or defend a smaller gap; do not swap it in. Reopen Millbrook at the start of
Lesson~1.7 --- the library \emph{could} assign, and that is exactly what 1.7 contrasts.
\end{teachernote}
\end{document}
